% translation-source: docs/user/latex/chapters/09-grading.tex; sha256: 9d1d80ab978ab206b1b2a0df2674872b7f4b71a1c8c6de39f41a0e42038bbe46
\chapter{Color grading}
\section{Choose a tonal region}
The buttons select shadows, midtones, highlights or the \emph{Global} wheel.
Each region retains its own settings independently. Unlike the mixer,
selection concerns a tonal region, not objects already having a particular color.

\section{Use the wheel}
Click the wheel and drag the selection. Angle selects hue; distance from
the center controls intensity. While dragging beyond the edge, interaction
continues at maximum intensity; there is no need to pick up the handle
again on each turn. The 0 to 100 \emph{Intensity} control also supports numeric
input. At zero, the region's color effect is neutral even if its hue is retained.

\section{Luminance, blending and balance}
Region luminance spans $-100$ to  $+100$. Blending, from 0 to 100, changes tonal
region overlap; balance, from$-100$ to  $+100$, shifts their balance.
The Global wheel applies an overall correction instead of selecting only
shadows, midtones or highlights.

\paragraph{Example.}
Give shadows a weak bluish tint and highlights a weak warm tint. Check
midtones and skin tones, then compare by holding \touche{Space} in the view.
Avoid pushing all three regions to maximum when assessing a natural rendition.

\paragraph{Keyboard.}
The intensity slider supports arrow keys, numeric input and double-click
reset as described for adjustments. Angular wheel control with arrow keys
is not currently declared.
